\section{Related Work}
\label{sec:related}

\para{Locate-then-edit methods.}
\citet{meng2022rome} introduced \rome, demonstrating that factual associations in GPT models correspond to localized, editable computations in mid-layer MLP modules.
\rome uses causal tracing to identify critical layers and applies a rank-one weight update to insert new key-value associations.
\citet{meng2023memit} extended this to \memit, which distributes updates across multiple MLP layers to support thousands of simultaneous edits.
\citet{fang2025alphaedit} proposed \alphaedit, which projects editing perturbations onto the null space of preserved-knowledge keys, providing a mathematical guarantee that preserved-knowledge outputs remain unchanged at the edited layer.
However, this guarantee is linear and applies only to the targeted FFN layer---non-linear interactions across layers may still produce side effects.
Our work tests whether these theoretical guarantees hold for computational knowledge, where the underlying representations may be more distributed than for factual associations.

\para{Meta-learning and memory-based approaches.}
\citet{mitchell2022mend} proposed \mend, a meta-learned hypernetwork that transforms fine-tuning gradients into targeted low-rank edits, though it degrades after approximately 10 sequential edits.
\citet{decao2021knowledgeeditor} trained a hypernetwork with constrained optimization to predict weight updates.
\citet{mitchell2022serac} introduced \serac, which stores edits in an external memory and routes relevant queries through a counterfactual model, avoiding weight modifications entirely.
These methods offer different trade-offs between edit capacity and model preservation, but none have been evaluated on arithmetic knowledge.

\para{Side effects and degradation from editing.}
A growing body of work documents unintended consequences of knowledge editing.
\citet{gu2024editing} showed that \rome, \memit, and \mend degrade general abilities including reasoning, natural language inference, and question answering---even when standard locality metrics appear satisfactory.
They proposed \rect, which regularizes the relative magnitude of weight changes to mitigate these effects.
\citet{li2024pitfalls} identified two failure modes: \emph{knowledge conflict}, where editing logically related facts magnifies inconsistencies, and \emph{knowledge distortion}, where parameter modifications irreversibly warp the model's internal knowledge structure.
\citet{gupta2024editing} demonstrated that sequential edits cause gradual forgetting followed by catastrophic collapse after 100--1000 edits, with even single edits shifting layer activation distributions.
Our work extends these findings to arithmetic knowledge, revealing a qualitatively different failure mode: systematic bias toward the injected answer rather than random degradation.

\para{Arithmetic and computational knowledge in LLMs.}
\citet{dai2022knowledge} identified ``knowledge neurons'' correlated with factual knowledge using gradient-based attribution, suggesting that some knowledge is neuron-localizable.
However, arithmetic may involve more complex computational circuits spanning multiple layers and attention heads.
\citet{li2023pmet} found that multi-head self-attention encodes general knowledge extraction patterns rather than specific facts, while FFN layers store the facts themselves.
Unlike prior work that focuses on entity-relation facts, we specifically target arithmetic---a form of \emph{computational} knowledge that may be stored in fundamentally different, more distributed representations.

\para{Low-rank adaptation.}
\citet{biderman2024lora} showed that \lora fine-tuning constrains models from diverging far from base weights, with an inverse relationship between fine-tuning performance and forgetting.
We use \lora as a baseline editing method, finding that its constrained parameter space does not prevent arithmetic side effects and in fact causes the most general language degradation among our tested methods.
